                  \sum_\xi|c(\xi)|^2\le B_0^2.
\end{equation}

\classification{\standardresult}{Finite Fourier inversion and Parseval; the source-specific centring is exact.}
\begin{proposition}[Meaning of zero line mean]\label{prop:mean-zero}
The identity \cref{eq:mean-zero} removes exactly the Fourier coefficient
$\xi=0$.  It supplies no pointwise decay for $c(\xi)$ when $\xi\ne0$ and does
not remove the diagonal frequencies $\xi=\eta\ne0$ in the two-copy expansion.
\end{proposition}

\begin{proof}
The first statement is the definition of $c(0)$.  A single character
$b^\perp(t)=\eord{Q}{\xi_0t}$ with $\xi_0\ne0$ is mean zero but has all of its
Fourier mass at $\xi_0$.  Moreover,
\begin{equation}\label{eq:two-copy-fourier}
 b(t)\overline{b(t+h)}
  =\sum_{\xi,\eta}c(\xi)\overline{c(\eta)}
       \eord{Q}{(\xi-\eta)t-\eta h}.
\end{equation}
The terms $\xi=\eta\ne0$ remain after $c(0)$ is deleted.
\end{proof}

Smoothness of the prime-box envelopes $V_{j,C}$ does not imply smoothness of
the sampled integer sequence $\beta_{6,C}(B_0+kt)$: changing $t$ changes the
factorisation condition.  No high-frequency Fourier decay may therefore be
inferred from the outer smooth weights alone.

\section{Closed local sectors before the rough-conductor pivot}\label{sec:closed-local}

The statements in this section belong to the frozen upstream analytic bank.
They are reproduced because they determine the exact residual, but the newest
independent audit did not re-prove each of them.  This provenance qualification
is part of the theorem statement rather than a footnote.

\classification{\newapplication}{Frozen upstream analytic bank; not Lean and not independently replayed in the newest rough-conductor audit.}
\begin{proposition}[Previously removed MHHC sectors]\label{prop:old-closed}
Subject to the original source hypotheses, the following sectors have been
removed from the MHHC source.
\begin{enumerate}[label=\textup{(\roman*)},leftmargin=2.7em]
\item The $h=0$ diagonal satisfies
      \begin{equation}\label{eq:diag-bound}
               \cO_{\mathrm{diag}}\ll X^{21/22+o(1)}.
      \end{equation}
\item The explicit KMT character modes satisfy
      \begin{equation}\label{eq:char-bound}
       \cO_{\mathrm{char}}
       \ll_B X\LL^{-B}+\LL^{O(1)}
          \left(K\sqrt{DA}+\frac{KD}{\sqrt{P_*}}\right)
       \ll X^{1-\delta_{\mathrm{char}}+o(1)},
      \end{equation}
      where
      \[
             \delta_{\mathrm{char}}
              =\min\left(\frac1{22},\epsstar\right).
      \]
\item The KMT-exceptional physical moduli contribute $O_B(X\LL^{-B})$.
\item For every fixed $B$, the near-diagonal range
      $0<|h|\le T\LL^{-B}$ contributes an arbitrarily large logarithmic
      saving after increasing the input exponent.
\item For every fixed $B$, the strata
      \[
       \max\{(d_t,h),(B_t,h)\}>\LL^B
      \]
      contribute an arbitrarily large logarithmic saving, together with a
      power-small endpoint error.
\end{enumerate}
\end{proposition}

\begin{proof}[Proof record and available calculations]
The diagonal estimate \cref{eq:diag-bound} is retained from the frozen bank;
its complete original source proof was not among the newest audit inputs and is
therefore not reconstructed here.

For the character modes, a nonprincipal character modulo the squarefree
five-prime $A$ has primitive conductor $f\mid A$ containing at least one
physical prime, so $f\ge P_*$.  The conductor-stratified multiplicative large
sieve gives the recorded bound
\[
 \sum_{\substack{A\asymp R\\\operatorname{cond}\chi_A\asymp F}}
 \frac1{\phi(A)}
 \left|\sum_{d\asymp D}\mu(d)\overline{\chi_A(d)}W_D(d/D)\right|
 \ll \LL^{O(1)}\sqrt{\frac{D(D+F^2)}F}.
\]
After the $k$-sum and dyadic summation in $F$, this yields the first inequality
in \cref{eq:char-bound}.  Since
$\log_X A\le5/11$, $\log_XK\le5/11$, and $D=X/K$,
\[
 \log_X\!\left(K\sqrt{DA}\right)
   =\frac{1+\log_XK+\log_XA}{2}\le\frac{21}{22},
\]
while $KD/\sqrt{P_*}\ll X^{1-\epsstar}$.

For exceptional $A$, the inherited application of
\cite{KMT2023} supplies
$\#\cE\ll AD^{-\eps^{200}}$.  With
\[
 \eps=\left(\frac{(B+C_0)\log\log X}{\log D}\right)^{1/200},
\]
one has $D^{-\eps^{200}}\le\LL^{-B-C_0}$.  The trivial mass per modulus is
$\ll KT\LL^{C_0}$ and $AKT\asymp X$, giving $O_B(X\LL^{-B})$.

For $H_0=T\LL^{-B}$, direct counting gives
\[
 \sum_{A,k}\sum_{0<|h|\le H_0}\sum_t|\cdots|
   \ll AKH_0T\LL^{O(1)}
   =AKT^2\LL^{-B+O(1)},
\]
which closes the near-diagonal after the finite ledger is included.

Finally,
\[
 (d_t,d_{t+h})=(d_t,h),\qquad
 (B_t,B_{t+h})=(B_t,h),
\]
because $(A,d_t)=(k,B_t)=1$.  For fixed $g\mid h$, the condition
$d_0+At\equiv0\pmod g$ places $t$ in at most one residue class modulo $g$.
Divisor switching then gives
\begin{equation}\label{eq:gcd-count}
 \#\{(h,t):|h|\le T,\ (d_t,h)>G\}
       \ll\frac{T^2}{G}+T\log T,
\end{equation}
and the same bound on the $B$ side.  Taking $G=\LL^B$ makes the first term
logarithmically small; since $T\ge X^{1/11-o(1)}$, the second is power-small
relative to the natural $T^2$ mass.
\end{proof}

\begin{remark}[Supersession of the old balanced-low estimate]
The previous KMT/endpoint-mesh argument gave only a relative gain
$\LL^{-1/500+o(1)}$ when
$K\ge R\LL^{-1/10000}$ and
$0<|\xi|\le\LL^{1/10000}$.  That estimate is no longer the controlling input.
The rough-modulus theorem proves arbitrary-logarithmic control there and in a
larger range.
\end{remark}


\section{The rough-modulus variance theorem}\label{sec:roughstatement}

The new arithmetic estimate is a consequence of the primitive multiplicative
large sieve, but the consequence is not obtained by simply replacing an
imprimitive character by its primitive inducer.  The entire coprimality
correction must be retained.  We restate the theorem in the notation used in
the proof.

Let $a_n\in\C$, $|a_n|\le1$, with $a_n=0$ outside $(D,2D]$.  For
$D\le y\le2D$, set
\[
 S_q(\chi;y)=\sum_{D<n\le y}a_n\chi(n).
\]
For $(a,q)=1$, define $E_q(a;y)$ as in
\cref{thm:roughvariance-main}.

\classification{\newderivation}{Derived from the standard primitive multiplicative large sieve; independently audited line by line; no claim of literature priority.}
\begin{theorem}[Rough-modulus variance, detailed form]\label{thm:roughvariance}
Under the hypotheses of \cref{thm:roughvariance-main},
\begin{equation}\label{eq:roughvariance}
 V(y):=\sum_{q\in\cA}\sum_{a\bmod q}^{*}|E_q(a;y)|^2
       \ll_j DR+\frac{D^2}{P}
\end{equation}
uniformly for $D\le y\le2D$.  The sequence $(a_n)$ must be common to all
$q$; modulus-dependent target information may not be absorbed into $a_n$.
\end{theorem}

The common-sequence restriction is exactly compatible with the physical
application: one takes $a_n=\mu(n)$, or $a_n=\mu(n)\1_{2\nmid n}$, while the
five-prime modulus, residue selected by $k$, and six-prime target remain outside
the source sequence.

\section{Complete proof of the rough-modulus variance theorem}\label{sec:roughproof}

\subsection{Character orthogonality}

For $(a,q)=1$, orthogonality of Dirichlet characters gives
\[
 \sum_{\substack{D<n\le y\\n\equiv a\pmod q}}a_n
 =\frac1{\phi(q)}\sum_{\chi\bmod q}\overline{\chi(a)}S_q(\chi;y).
\]
The principal-character term is precisely
$\phi(q)^{-1}\sum_{D<n\le y,(n,q)=1}a_n$.  Consequently,
\begin{equation}\label{eq:E-character}
 E_q(a;y)=\frac1{\phi(q)}
     \sum_{\substack{\chi\bmod q\\\chi\ne\chi_0}}
        \overline{\chi(a)}S_q(\chi;y),
\end{equation}
and a second application of orthogonality over reduced residues yields the
exact identity
\begin{equation}\label{eq:V-character}
 V(y)=\sum_{q\in\cA}\frac1{\phi(q)}
       \sum_{\substack{\chi\bmod q\\\chi\ne\chi_0}}
         |S_q(\chi;y)|^2.
\end{equation}
There is no contribution from non-coprime $n$ on either side.

\subsection{Exact imprimitive descent}

Fix a nonprincipal character $\chi\bmod q$, and let $\chi^*$ be the unique
primitive character inducing it, of conductor $f>1$.  Write $q=fr$.  Because
$q$ is squarefree, $(f,r)=1$, and every prime dividing $f$ is at least $P$;
therefore
\begin{equation}\label{eq:f-lower}
                         f\ge P.
\end{equation}
For every integer $n$,
\[
              \chi(n)=\chi^*(n)\1_{(n,r)=1}.
\]
M\"obius inversion of the coprimality indicator gives
\begin{align}
 S_q(\chi;y)
 &=\sum_{D<n\le y}a_n\chi^*(n)
          \sum_{e\mid(n,r)}\mu(e)\notag\\
 &=\sum_{e\mid r}\mu(e)\chi^*(e)
       \sum_{D/e<m\le y/e}a_{em}\chi^*(m).
       \label{eq:imprimitive-exact}
\end{align}
The factor is $\chi^*(e)$, not its conjugate.  Since $e\mid r$ and
$(f,r)=1$, one has $(e,f)=1$ automatically.

As $r$ has at most $j$ prime factors,
$\tau(r)\le2^j$.  Cauchy--Schwarz in \cref{eq:imprimitive-exact} gives
\begin{equation}\label{eq:imprimitive-cauchy}
 |S_q(\chi;y)|^2
 \le2^j\sum_{e\mid r}
       \left|\sum_{D/e<m\le y/e}a_{em}\chi^*(m)\right|^2.
\end{equation}
Moreover,
\begin{equation}\label{eq:phi-lower}
        \phi(q)=q\prod_{p\mid q}(1-p^{-1})\ge q2^{-j}\ge R2^{-j},
\end{equation}
so substitution into \cref{eq:V-character} yields
\begin{equation}\label{eq:V-preblock}
 V(y)\le\frac{4^j}{R}
 \sum_{f>1}\sum_{\chi^*\, (\mathrm{mod}\ f)}^{\mathrm{prim}}
 \sum_{\substack{r:\,fr\in\cA}}
 \sum_{e\mid r}
 \left|\sum_{D/e<m\le y/e}a_{em}\chi^*(m)\right|^2.
\end{equation}
All omitted restrictions in later steps are removed only from nonnegative sums.

\subsection{The conductor block and the $fe>R$ edge}

Restrict $f$ to a dyadic interval $F\le f<2F$.  For fixed $f$ and $e$, the
number of integers $r$ satisfying $e\mid r$ and $R\le fr<2R$ is at most
\begin{equation}\label{eq:r-count}
                         \frac{2R}{fe}.
\end{equation}
Here is the endpoint check that was implicit in the first draft.  Write
$r=em$.  If $fe\le R$, the interval for $m$ has length $R/(fe)$, so the count
is at most $R/(fe)+1\le2R/(fe)$.  If $R<fe<2R$, at most the single integer
$m=1$ can occur, and $1\le2R/(fe)$.  If $fe\ge2R$, there is no permissible
$m$.  Thus \cref{eq:r-count} is uniform and no stray ``$+1$'' survives.

Using $f\ge F$, extending the $e$-sum, and relaxing $f<2F$ to $f\le2F$ gives
\begin{equation}\label{eq:VF}
 V_F(y)\ll 4^j\sum_{e\ge1}\frac1{Fe}
 \sum_{f\le2F}\sum_{\chi^*\, (\mathrm{mod}\ f)}^{\mathrm{prim}}
 \left|\sum_{D/e<m\le y/e}a_{em}\chi^*(m)\right|^2.
\end{equation}
This is the point at which the factor $1/R$ from \cref{eq:V-preblock} is
cancelled by the number of possible complementary moduli $r$.
